%=======================================================================
%  NeuroPresence -- Proof-of-Concept UI Prototype
%  Complete Feature and Implementation Specification
%
%  Compiles with pdfLaTeX. No non-standard packages, no external assets.
%  Deliberately ASCII-only in the source so it builds anywhere.
%=======================================================================
\documentclass[11pt,a4paper]{report}

\usepackage[T1]{fontenc}
\usepackage[utf8]{inputenc}
\usepackage{lmodern}
\usepackage[margin=2.4cm,top=2.6cm,bottom=2.6cm]{geometry}
\usepackage[table]{xcolor}
\usepackage{booktabs}
\usepackage{longtable}
\usepackage{array}
\usepackage{tabularx}
\usepackage{enumitem}
\usepackage{titlesec}
\usepackage{fancyhdr}
\usepackage{amsmath}
\usepackage{amssymb}
\usepackage{graphicx}
\usepackage{listings}
\usepackage[hidelinks]{hyperref}

%----------------------------------------------------------------------
%  Palette and styling
%----------------------------------------------------------------------
\definecolor{npPrimary}{HTML}{2563EB}
\definecolor{npAccent}{HTML}{7C3AED}
\definecolor{npSuccess}{HTML}{16A34A}
\definecolor{npWarning}{HTML}{D97706}
\definecolor{npDanger}{HTML}{DC2626}
\definecolor{npInk}{HTML}{0F172A}
\definecolor{npMuted}{HTML}{64748B}
\definecolor{npRule}{HTML}{E2E8F0}
\definecolor{npSurface}{HTML}{F3F5F9}

\hypersetup{
  colorlinks=false,
  pdftitle={NeuroPresence Prototype -- Feature and Implementation Specification},
  pdfauthor={Zain Shahid, Muhammad Talha Arshad, Sana Ullah Farooqi},
  pdfsubject={FYP-1 proposal defense: proof-of-concept UI prototype},
}

\titleformat{\chapter}[display]
  {\normalfont\huge\bfseries\color{npInk}}
  {\color{npPrimary}\normalsize\MakeUppercase{Chapter \thechapter}}
  {6pt}{\Huge}
\titlespacing*{\chapter}{0pt}{-24pt}{24pt}

\titleformat{\section}{\normalfont\large\bfseries\color{npInk}}{\thesection}{0.6em}{}
\titleformat{\subsection}{\normalfont\normalsize\bfseries\color{npInk}}{\thesubsection}{0.6em}{}

\pagestyle{fancy}
\fancyhf{}
\renewcommand{\headrulewidth}{0.4pt}
\fancyhead[L]{\small\color{npMuted}NeuroPresence -- Prototype Specification}
\fancyhead[R]{\small\color{npMuted}\thepage}

\setlist[itemize]{leftmargin=1.35em,itemsep=2pt,topsep=3pt}
\setlist[enumerate]{leftmargin=1.6em,itemsep=2pt,topsep=3pt}

\newcolumntype{L}[1]{>{\raggedright\arraybackslash}p{#1}}
\newcommand{\code}[1]{\texttt{\small #1}}

% A boxed caution note without needing tcolorbox.
\newenvironment{npnote}[1]
  {\par\medskip\noindent\begin{tabularx}{\linewidth}{|>{\hspace{4pt}}X<{\hspace{4pt}}|}
   \hline\rule{0pt}{3ex}\textbf{#1}\par\medskip}
  {\par\medskip\hline\end{tabularx}\par\medskip}

\lstset{
  basicstyle=\ttfamily\scriptsize,
  frame=single,
  rulecolor=\color{npRule},
  backgroundcolor=\color{npSurface},
  breaklines=true,
  columns=fullflexible,
  keepspaces=true,
  showstringspaces=false,
  xleftmargin=0pt,
  framexleftmargin=3pt,
}

%=======================================================================
\begin{document}

%----------------------------------------------------------------------
%  Title page
%----------------------------------------------------------------------
\begin{titlepage}
\centering
\vspace*{2.2cm}

{\color{npPrimary}\rule{\linewidth}{2pt}}\\[10pt]
{\Huge\bfseries NeuroPresence\par}
\vspace{6pt}
{\Large Proof-of-Concept UI Prototype\par}
\vspace{4pt}
{\large\itshape Complete Feature and Implementation Specification\par}
{\color{npPrimary}\rule{\linewidth}{2pt}}\\[26pt]

{\large Composed presence for online meetings,\\ using the speaker's own consented likeness\par}
\vspace{28pt}

\begin{tabular}{rl}
\toprule
\textbf{Project} & NeuroPresence \\
\textbf{Stage} & FYP-1 proposal defense \\
\textbf{Artifact} & High-fidelity interactive UI prototype \\
\textbf{University} & FAST-NUCES Islamabad \\
\textbf{Supervisor} & Muhammad Aamir Gulzar \\
\bottomrule
\end{tabular}

\vspace{20pt}
{\bfseries Team}\\[4pt]
\begin{tabular}{ll}
\toprule
Zain Shahid & i232582 \\
Muhammad Talha Arshad & i232548 \\
Sana Ullah Farooqi & i232594 \\
\bottomrule
\end{tabular}

\vfill

\begin{tabularx}{0.94\linewidth}{|>{\hspace{5pt}}X<{\hspace{5pt}}|}
\hline
\rule{0pt}{3.2ex}%
{\bfseries\color{npDanger} Scope statement.}
This document describes a \emph{user-interface prototype}. The reenactment
pipeline is \textbf{not implemented} in this build. There is no neural
network, no LivePortrait, no PyTorch, ONNX or TensorRT, no face tracking, no
landmark detection, no identity matching, no virtual-camera device, no
backend, no database and no network call. Every metric, gate decision and
render produced on screen originates in the simulation module described in
Chapter~\ref{ch:engine}. The measured figures quoted in Chapter~\ref{ch:numbers}
come from a separate benchmark run on the team's own hardware and are the
only non-simulated numbers in the system.\par\medskip
\hline
\end{tabularx}

\vspace{14pt}
{\small\color{npMuted}Document version 1.0}
\end{titlepage}

\tableofcontents
\cleardoublepage

%=======================================================================
\chapter{Introduction}
%=======================================================================

\section{The product concept}

NeuroPresence lets a person appear composed and professional in an online
meeting without having to be camera-ready at that moment. The user records one
short, presentable \emph{source clip} of themselves in advance. During a
meeting their live webcam supplies a \emph{driving signal} -- head pose, lip
motion and expression. A reenactment model warps the source clip to follow that
motion, so the outgoing video looks like the user, composed, saying what they
are saying right now. The result is written to a \emph{virtual camera}, so any
meeting application reads it as an ordinary webcam with no plugin.

Two safeguards are integral rather than optional. A \emph{consent gate}
restricts animation to the enrolled user's own likeness, and an optional
\emph{synthetic-media disclosure watermark} tells other participants the feed
is reenacted.

\section{Two operating modes}

\begin{description}[leftmargin=0pt,style=nextline]
  \item[Real-time mode (primary).] Driven by the live webcam and bounded by a
    latency budget, for use during live meetings.
  \item[Offline mode.] The same feature set driven by a pre-recorded video the
    user uploads. With the time constraint removed it can run at higher
    fidelity; it serves asynchronous recording and acts as a quality reference
    against which real-time output is judged.
\end{description}

\section{What this artifact is}

The evaluating committee defines a proof of concept as \emph{a high-fidelity
functional prototype: an interactive, highly detailed model that closely
mirrors the visual design and behaviour of a finished product; a skeleton of
the finished product without actual implementation of the functionality behind
the interface.}

This prototype is built to that definition. The interface is real, complete and
interactive across eight screens. The functionality behind it is deliberately
absent, and every screen states so.

\section{Purpose of this document}

This is a complete inventory of the prototype as built: every screen, every
control, every state, every component, every simulated behaviour and every
constant. It is intended to be read alongside the running application, and to
serve as the reference from which the implemented system is later specified.

%=======================================================================
\chapter{Technology and Project Structure}
%=======================================================================

\section{Stack}

\begin{longtable}{L{4.6cm} L{3.2cm} L{7.2cm}}
\toprule
\textbf{Component} & \textbf{Version} & \textbf{Role} \\
\midrule
\endhead
React & 18.3.1 & UI runtime; hooks and context only, no external state library \\
React DOM & 18.3.1 & Rendering, including portals for overlays \\
TypeScript & 5.7.2 & Static typing across the whole source tree \\
Vite & 6.0.7 & Dev server and static production build \\
Tailwind CSS & 3.4.17 & Utility styling over a custom token layer \\
Framer Motion & 11.15.0 & Transitions, layout animation, scroll effects \\
lucide-react & 0.469.0 & Icon set \\
PostCSS / Autoprefixer & 8.4.49 / 10.4.20 & CSS pipeline \\
\bottomrule
\end{longtable}
\noindent\textit{\small Runtime and build dependencies. No machine-learning package appears
anywhere in the dependency tree.}\par\medskip

\section{Commands}

\begin{longtable}{L{5.2cm} L{9.6cm}}
\toprule
\textbf{Command} & \textbf{Effect} \\
\midrule
\endhead
\code{npm install} & Install dependencies \\
\code{npm run dev} & Development server on \code{http://localhost:5173} \\
\code{npm run build} & Type-check, then emit static files to \code{dist/} \\
\code{npm run preview} & Serve the production build locally \\
\code{npm run typecheck} & Type-check without emitting \\
\bottomrule
\end{longtable}

\section{Source tree}

The source is roughly 8{,}830 lines of TypeScript and TSX across 39 files.

\begin{longtable}{L{5.4cm} L{9.4cm}}
\toprule
\textbf{Path} & \textbf{Contents} \\
\midrule
\endhead
\code{src/main.tsx} & Entry point; mounts the provider and the app \\
\code{src/App.tsx} & Shell composition and screen switch \\
\code{src/index.css} & Design tokens, base layer, component classes \\
\code{src/app/} & \code{routes.ts}, \code{Sidebar.tsx}, \code{TopBar.tsx} \\
\code{src/screens/} & Eight screen components \\
\code{src/components/} & Seventeen UI and domain components \\
\code{src/mock/} & Simulation: engine, sampler, seed data, constants, tour, evidence, data checks, types \\
\bottomrule
\end{longtable}

Everything that stands in for the absent pipeline lives under
\code{src/mock/}. This separation is deliberate: the boundary between what is
real interface and what is simulated behaviour is a directory boundary, and can
be audited as one.

%=======================================================================
\chapter{Design System}
%=======================================================================

\section{Colour tokens}

Colours are declared as RGB channel triples on the document root so that every
Tailwind utility can compose an alpha value against them
(\code{bg-success/10}, \code{border-danger/40}). Resolved \code{--token} forms
remain available to hand-written CSS.

\begin{longtable}{L{3.2cm} L{2.6cm} L{2.6cm} L{5.6cm}}
\toprule
\textbf{Token} & \textbf{Light} & \textbf{Dark} & \textbf{Use} \\
\midrule
\endhead
\code{--bg} & \code{\#F9FAFC} & \code{\#0B0F14} & Page ground \\
\code{--surface} & \code{\#FFFFFF} & \code{\#141A22} & Cards and panels \\
\code{--surface-2} & \code{\#F3F5F9} & \code{\#1C2530} & Inputs, raised elements \\
\code{--border} & \code{\#E2E8F0} & \code{\#26313D} & Hairlines, card borders \\
\code{--text} & \code{\#0F172A} & \code{\#E6EDF3} & Primary text \\
\code{--text-muted} & \code{\#64748B} & \code{\#8B98A5} & Secondary text \\
\code{--primary} & \code{\#2563EB} & \code{\#3B82F6} & Primary action, active state \\
\code{--primary-hover} & \code{\#1D4ED8} & \code{\#2563EB} & Hover \\
\code{--success} & \code{\#16A34A} & \code{\#22C55E} & Targets met, live OK \\
\code{--warning} & \code{\#D97706} & \code{\#F59E0B} & Caution, baseline figures \\
\code{--danger} & \code{\#DC2626} & \code{\#EF4444} & Errors, stop, gate blocked \\
\code{--accent} & \code{\#7C3AED} & \code{\#8B5CF6} & Highlights, disclosure chip \\
\bottomrule
\end{longtable}
\noindent\textit{\small Colour tokens. The dark palette is the one specified in the build
brief, reproduced exactly.}\par\medskip

Light is the product default. Dark is one click away in Settings and via a
toggle in the landing-page header. Elevation reads through shadow in light and
through borders in dark, so four shadow tokens
(\code{--shadow-card}, \code{--shadow-raised}, \code{--shadow-hover},
\code{--shadow-hero}) are redefined per theme rather than shared.

Two further per-theme tokens support the landing page: \code{--grid-line} for
the faint engineering grid behind hero sections, and \code{--glow-a} /
\code{--glow-b} for the drifting colour blooms.

\section{Typography}

\begin{longtable}{L{4.4cm} L{10.4cm}}
\toprule
\textbf{Role} & \textbf{Specification} \\
\midrule
\endhead
Interface face & Inter, falling back to system sans stacks \\
Numeric face & JetBrains Mono, falling back to \code{ui-monospace} \\
Page title & 24\,px / weight 600 / tight leading \\
Section heading & 16\,px / weight 600 \\
Body & 14\,px / weight 400 / 1.55 leading \\
Metric value & 28\,px / weight 600 / monospace, tabular figures \\
Metric label & 11--12\,px / weight 600 / uppercase / 0.09\,em tracking \\
\bottomrule
\end{longtable}

Metric values use tabular numerals so a changing figure never reflows its
tile. Headings carry $-0.02$\,em tracking.

\section{Spacing, radius, layout}

An 8\,px spacing grid throughout. Card padding is 20\,px by default. Corner
radii: cards 14\,px, controls 9\,px, chips 6\,px, hero surfaces 20\,px. The
landing page is constrained to 1180\,px; the application shell to 1480\,px. The
design reference is 1440$\times$900, and the layout remains usable down to
480\,px.

\section{Motion}

State changes run at 150--250\,ms with an ease-out curve of
\code{cubic-bezier(0.16, 1, 0.3, 1)}. Nine named keyframe animations are
defined: \code{pulse-dot}, \code{ring-sweep}, \code{fade-up}, \code{shimmer},
\code{gradient-pan}, \code{float}, \code{float-slow}, \code{marquee} and
\code{dash-flow}.

A global \code{prefers-reduced-motion} rule collapses every animation and
transition to 0.001\,ms and disables smooth scrolling. Individual components
additionally branch on \code{useReducedMotion()} so that scroll-triggered
reveals, the count-up numerals, the tilt effect and the pipeline flow dots
resolve to their final state instead of animating.

\section{Component classes}

Six semantic classes are defined in the CSS component layer so repeated
patterns stay consistent: \code{.np-card}, \code{.np-chip}, \code{.np-label},
\code{.np-metric}, \code{.np-section-heading}, \code{.np-page-title},
\code{.np-divider}. The landing page adds \code{.np-gradient-text},
\code{.np-grid-bg}, \code{.np-eyebrow} and \code{.np-rule}.

%=======================================================================
\chapter{Application Shell and Navigation}
%=======================================================================

\section{Shell composition}

The landing page and the enrollment flow each own the whole viewport. Every
other screen renders inside a persistent shell:

\begin{itemize}
  \item \textbf{Left sidebar.} Logo and wordmark (returns to the overview);
    six navigation items; a guided-demo launcher; a project-overview link; the
    enrolled-user chip with an avatar, name and status dot; and a permanent
    \emph{Prototype build -- simulated data} badge. The sidebar is 252\,px wide
    and collapses to a 72\,px icon rail below 1280\,px.
  \item \textbf{Top bar.} Screen title and subtitle; the Real-time / Offline
    mode switch; a consent chip; a disclosure chip; a session-state chip
    (Idle / Warming / LIVE); and the virtual-camera output name.
  \item \textbf{Main region.} The active screen, followed by a footer carrying
    the prototype disclosure.
\end{itemize}

The active navigation item is marked with a shared-layout highlight that
springs between items rather than cutting.

\section{Routing}

Routes are hash-based rather than history paths, because the build ships as
static files and a hash URL resolves on a hard refresh without server rewrite
rules.

\begin{longtable}{L{4.6cm} L{4.6cm} L{5.6cm}}
\toprule
\textbf{Route} & \textbf{Screen} & \textbf{Document title} \\
\midrule
\endhead
\code{\#/} & Landing / overview & NeuroPresence -- Composed presence \ldots \\
\code{\#/enroll} & Onboarding & Get started \\
\code{\#/console} & Console & Console \\
\code{\#/source-clips} & Source Clips & Source Clips \\
\code{\#/offline-studio} & Offline Studio & Offline Studio \\
\code{\#/devices} & Devices \& Output & Devices \& Output \\
\code{\#/feasibility} & Feasibility evidence & Feasibility Evidence \\
\code{\#/settings} & Settings & Settings \\
\bottomrule
\end{longtable}

Behaviour guaranteed by the router:

\begin{itemize}
  \item The document title follows the route.
  \item Browser back and forward move between screens.
  \item A hard refresh lands on the same screen.
  \item A pasted deep link opens that screen directly. The initial screen is
    seeded from the URL on the first render, so the screen-to-hash effect can
    never fire with a stale value and overwrite the requested route.
  \item An unknown hash falls back to the landing page.
  \item Landing-page section links scroll programmatically and never write an
    anchor into the URL, so \code{\#/\ldots} always denotes a screen.
  \item Deep-linking past enrollment adopts the demo identity, so the shell is
    never half-enrolled.
\end{itemize}

%=======================================================================
\chapter{State Model: the Simulation Engine}
\label{ch:engine}
%=======================================================================

All simulated state lives in a single React context provider. There is no
persistence of any kind -- no \code{localStorage}, no \code{sessionStorage}, no
backend -- so every reload starts from a clean, predictable state. This is a
deliberate property: a demonstration cannot be spoiled by residue from a
previous run.

\section{Public interface}

\begin{longtable}{L{4.9cm} L{4.0cm} L{6.0cm}}
\toprule
\textbf{Member} & \textbf{Type} & \textbf{Meaning} \\
\midrule
\endhead
\multicolumn{3}{l}{\textit{Navigation}}\\
\code{screen} & \code{Screen} & Active screen \\
\code{navigate} & \code{(s) => void} & Change screen \\
\code{mode} & \code{'realtime' | 'offline'} & Operating mode \\
\code{setMode} & \code{(m) => void} & Switch mode; routes to the matching screen \\
\addlinespace
\multicolumn{3}{l}{\textit{Identity}}\\
\code{enrolledUser} & \code{EnrolledUser | null} & Name, label identifier, timestamp \\
\code{enrollUser} & \code{(name) => void} & Record an enrollment \\
\code{onboardingDone} & \code{boolean} & Enrollment flow completed \\
\code{completeOnboarding} & \code{() => void} & Finish and open the Console \\
\code{restartEnrollment} & \code{() => void} & Re-enroll from Settings \\
\addlinespace
\multicolumn{3}{l}{\textit{Source clips}}\\
\code{clips} & \code{SourceClip[]} & Library, seeds plus uploads \\
\code{activeClip} & \code{SourceClip} & The clip that gets animated \\
\code{setActiveClip} & \code{(id) => void} & Change the active clip \\
\code{addClip} & \code{(file) => SourceClip} & Register an uploaded file \\
\addlinespace
\multicolumn{3}{l}{\textit{Session}}\\
\code{session} & \code{'idle'|'warming'|'live'} & Session state \\
\code{warmupProgress} & \code{number} & 0--100 during warm-up \\
\code{startSession} & \code{() => void} & Begin; no-op if already running \\
\code{stopSession} & \code{() => void} & End and reset the clock \\
\code{sessionSeconds} & \code{number} & Elapsed live seconds \\
\addlinespace
\multicolumn{3}{l}{\textit{Consent gate}}\\
\code{gateEnabled} & \code{boolean} & Policy switch \\
\code{gateState} & \code{'verified'|'disabled'|'blocked'} & Derived presentation state \\
\code{simulateNonEnrolledFace} & \code{() => void} & Demonstration: block for 4\,s \\
\code{blockedSecondsLeft} & \code{number} & Recovery countdown \\
\addlinespace
\multicolumn{3}{l}{\textit{Disclosure}}\\
\code{watermark} & \code{boolean} & Overlay on or off \\
\addlinespace
\multicolumn{3}{l}{\textit{Telemetry}}\\
\code{metricsMode} & \code{'target'|'baseline'} & Which figures are shown \\
\code{metrics} & \code{MetricSample | null} & Latest sample \\
\code{history} & \code{MetricSample[]} & Rolling 30-sample window \\
\addlinespace
\multicolumn{3}{l}{\textit{Devices}}\\
\code{inputDevice} & \code{string} & Selected camera label \\
\addlinespace
\multicolumn{3}{l}{\textit{Offline studio}}\\
\code{offline} & \code{OfflineState} & Upload, checks, settings, render \\
\code{setDrivingVideo} & \code{(file) => void} & Accept an upload \\
\code{loadDemoDrivingVideo} & \code{() => void} & Stand-in recording for the tour \\
\code{startRender} / \code{resetRender} & \code{() => void} & Render lifecycle \\
\addlinespace
\multicolumn{3}{l}{\textit{Appearance and demo}}\\
\code{theme} & \code{'dark'|'light'} & Applied as a root data attribute \\
\code{tourActive}, \code{tourIndex}, \code{tourPlaying} & --- & Guided-demo state \\
\code{startTour} \ldots \code{tourGoTo} & --- & Guided-demo transport \\
\bottomrule
\end{longtable}
\noindent\textit{\small The complete engine surface. Forty-eight members in total.}\par\medskip

\section{Session lifecycle}

\[
\text{idle} \xrightarrow{\ \text{Start}\ } \text{warming}
\xrightarrow{\ 1500\,\text{ms}\ } \text{live}
\xrightarrow{\ \text{Stop}\ } \text{idle}
\]

Warm-up progress is sampled every 50\,ms and reported as a percentage. While
live, a clock ticks once per second and the telemetry sampler emits twice per
second. Calling \code{startSession} on a session that is already warming or
live is a no-op, so a repeated press cannot restart the warm-up.

A blocked consent gate tears down a running session immediately, whichever
state it is in, and disables the start control until the gate recovers.

\section{Consent gate}

The gate state is derived rather than stored:

\begin{center}
\begin{tabular}{ll}
\toprule
\textbf{Condition} & \textbf{State} \\
\midrule
Policy switch off & \code{disabled} (amber) \\
Blocked timer active & \code{blocked} (red) \\
Otherwise & \code{verified} (green) \\
\bottomrule
\end{tabular}
\end{center}

The demonstration control sets a four-second block. A 200\,ms interval drives a
visible countdown, and the state recovers on its own.

\section{Timing constants}

\begin{center}
\begin{tabular}{lll}
\toprule
\textbf{Constant} & \textbf{Value} & \textbf{Governs} \\
\midrule
\code{WARMUP\_MS} & 1500 & Session warm-up \\
\code{BLOCKED\_DEMO\_MS} & 4000 & Consent-gate block \\
\code{METRICS\_INTERVAL\_MS} & 500 & Telemetry sampling \\
\code{RENDER\_MS} & 5000 & Offline render \\
\code{ANALYSIS\_MS} & 2000 & Enrollment analysis \\
Input check & 1400 & Offline input-quality check \\
\code{HISTORY\_LENGTH} & 30 & Sparkline window \\
\bottomrule
\end{tabular}
\end{center}

\section{Telemetry simulation}

Samples are produced by a mean-reverting random walk bounded to the configured
envelope. For a previous value $x_{t-1}$ in a range $[\ell, h]$ with midpoint
$m = (\ell + h)/2$ and span $s = h - \ell$:

\[
x_t = \operatorname{clamp}\Big(x_{t-1} + \big(m - x_{t-1}\big)\,\alpha
      + \mathcal{N}_{\pm 2.5}\,s\,\beta,\ \ \ell,\ h\Big)
\]

where $\mathcal{N}_{\pm 2.5}$ is a Box--Muller normal deviate clamped to
$\pm 2.5$ standard deviations so a single sample can never fly off. Pull
$\alpha$ and noise $\beta$ default to $0.12$ and $0.16$; GPU memory uses
$0.05$ and $0.07$ because allocators do not thrash frame to frame, and
identity similarity uses $0.09$ and $0.12$.

The result wanders and occasionally ticks, which is what real telemetry looks
like, rather than jittering uniformly. The function is pure: the next state is
a function of the previous state and the mode alone.

\begin{longtable}{L{3.4cm} L{4.4cm} L{4.4cm}}
\toprule
\textbf{Metric} & \textbf{Target envelope} & \textbf{Baseline envelope} \\
\midrule
\endhead
Latency & 38--45\,ms & 127.1--133.1\,ms \\
Throughput & 22--25\,FPS & 7.4--8.0\,FPS \\
GPU memory & 3.5--4.5\,GB & 1.53--1.61\,GB \\
Identity (CSIM) & 0.83--0.88 & reported as \emph{n/a (baseline)} \\
\bottomrule
\end{longtable}
\noindent\textit{\small Simulation envelopes. Baseline envelopes are the measured constants
with small jitter; the baseline benchmark produced no identity figure, so none
is invented.}\par\medskip

The sampler runs only while a session is live. On stop, values freeze at the
last sample rather than resetting, so the final reading remains readable.

\section{Seed data}

\begin{longtable}{L{4.4cm} L{2.0cm} L{2.4cm} L{5.2cm}}
\toprule
\textbf{Clip} & \textbf{Duration} & \textbf{Resolution} & \textbf{Scene} \\
\midrule
\endhead
Office --- Navy Blazer & 18\,s & 512$\times$512 & office \\
Home Study & 24\,s & 512$\times$512 & study \\
Plain Wall & 12\,s & 512$\times$512 & plain \\
\bottomrule
\end{longtable}

Each seed clip carries a six-colour palette (background gradient, key light,
garment, skin, hair) used by the procedural renderer. No video files are
shipped: the repository stays small, and nothing on screen can be mistaken for
model output. Two further palettes cycle for user-added clips.

Additional seed collections: three input devices (FaceTime HD Camera, USB
Webcam, Integrated Camera 1080p); three meeting clients (Zoom, Google Meet,
Microsoft Teams) with brand tints; four render stage captions (Aligning
face, Transferring motion, Compositing, Finalizing); and four input-quality
checks (Face detected, Front-facing, Adequate lighting, Single face).

%=======================================================================
\chapter{Screens}
%=======================================================================

\section{Landing / project overview}

The entry screen. A single scrolling page that explains the whole project,
owning the full viewport with its own sticky navigation.

\begin{longtable}{L{3.6cm} L{11.2cm}}
\toprule
\textbf{Section} & \textbf{Contents} \\
\midrule
\endhead
Navigation & Wordmark, five section links with a spring-animated active
  indicator driven by a scroll spy, a theme toggle, and a launch button.
  Collapses to a menu below 1024\,px. \\
Hero & Word-by-word revealed headline with a gradient accent phrase, the
  approved tagline, three calls to action (launch, watch the guided tour, jump
  to the measured numbers), a drifting aurora backdrop, an engineering grid,
  and a product shot that parallaxes on scroll. \\
Stat strip & Headline figures drawn from the measured run and the targets. \\
Problem & The presentation tax that the product removes. \\
How it works & The pipeline as a sequence of stages with animated connectors. \\
Two modes & Real-time and offline compared side by side. \\
Safeguards & The consent gate and the disclosure watermark, with the three
  gate states shown. \\
Proof of concept & Measured baseline against engineering target as scaled
  bars with count-up numerals; full measurement and target tables; a link into
  the Feasibility screen; and an explicit note on where honesty sits in the
  build. \\
Capabilities & The functional-requirement map. \\
Integration & Meeting-client support and the desktop-only limitation. \\
Team & Members, roll numbers, supervisor and university. \\
Final call to action & Entry into the prototype. \\
Footer & Prototype disclosure. \\
\bottomrule
\end{longtable}

A reading-progress bar is pinned to the top of the viewport.

\section{Onboarding / identity enrollment}

A centred card with a three-step progress indicator.

\begin{enumerate}
  \item \textbf{Welcome.} Product name, tagline, a \emph{Get started} action,
    the enrollment note, and a \emph{Skip -- use demo identity} link that
    jumps straight to the Console.
  \item \textbf{Capture identity.} A circular framed preview, a name field
    defaulting to the first team member, and a \emph{Capture} action. Capture
    runs a two-second progress indicator captioned \emph{Analyzing facial
    embedding}, then resolves to a green \emph{Identity enrolled} confirmation
    showing the name and a label identifier. The step is a timer; no embedding
    is computed. If camera permission has already been granted the real webcam
    is shown; otherwise an avatar with an animated scanning ring stands in, and
    an opt-in button offers the camera without ever prompting unbidden.
  \item \textbf{Choose a source clip.} A grid of the available clips with
    selection state, and a \emph{Finish} action that lands on the Console.
\end{enumerate}

Re-enrolling from Settings enters at step two.

\section{Console}

The centrepiece and the application home. A two-column layout.

\subsection*{Left column}

\begin{itemize}
  \item \textbf{Virtual Camera Output frame.} A 16:9 surface showing the active
    source clip as a stand-in for reenacted output, with a resolution chip, a
    frame label, and a live indicator with a pulsing dot. Four overlay states:
    an idle veil, a warm-up progress gradient, a red blocked veil when the gate
    refuses, and the disclosure watermark when enabled.
  \item \textbf{Session controls.} Start or Stop, a warming state, an elapsed
    clock in monospace, and a signal-path readout reading
    \emph{Driving: Webcam} $\rightarrow$ \emph{NeuroPresence Camera}.
  \item \textbf{Reenactment pipeline strip.} Five stages -- Webcam, Motion,
    Consent gate, Reenactment, Virtual camera -- each labelled with the
    functional requirement it satisfies. While live, particles flow along the
    connectors and stages pulse; when the gate blocks, the strip halts at the
    gate and everything downstream dims. A per-frame latency chip and a
    \emph{simulated} badge sit in the header.
  \item \textbf{Active source clip.} Name, duration and resolution, with a
    \emph{Change} action, and a note stating that the preview is a stand-in.
  \item \textbf{Driving signal.} An optional webcam passthrough. Permission is
    requested only on explicit press. Frames are displayed and never analysed.
\end{itemize}

\subsection*{Right column}

\begin{itemize}
  \item \textbf{Live telemetry.} Described in Section~\ref{sec:metrics}.
  \item \textbf{Consent gate card.} State chip, avatar with a state-coloured
    ring, the approved message for the current state, a policy toggle, and the
    non-enrolled-face demonstration with its countdown.
  \item \textbf{Disclosure watermark card.} A toggle and a live preview
    thumbnail of the overlay as it appears in the output.
\end{itemize}

\subsection{Live telemetry panel}
\label{sec:metrics}

Four tiles, a mode switch, a history sparkline and a permanent
\emph{simulated} badge.

\begin{longtable}{L{3.4cm} L{2.6cm} L{8.6cm}}
\toprule
\textbf{Tile} & \textbf{Unit} & \textbf{Status thresholds} \\
\midrule
\endhead
Latency & ms & Green $\leq 42$; amber $\leq 150$; red above \\
Throughput & FPS & Green $\geq 24$; amber $\geq 12$; red below \\
GPU memory & GB of 8.0 & Green $\leq 8$; red above \\
Identity (CSIM) & 0--1 & Green $\geq 0.80$; amber below; \emph{n/a} in baseline \\
\bottomrule
\end{longtable}

Values ease toward each new sample rather than jumping, using a
requestAnimationFrame interpolator; a change large enough to be a mode switch
rather than a new sample lands immediately instead of sliding through
meaningless intermediate figures. Tiles use tabular numerals so nothing
reflows.

The sparkline draws the rolling 30-sample latency history as a Catmull--Rom
spline emitted as cubic Beziers, with a gradient area fill, a soft under-glow,
a dashed target threshold line and a breathing marker on the leading sample.

The \textbf{Target / Baseline} switch is the panel's most important control. In
Target it shows the finished-product operating point; in Baseline it clamps to
the measured proof-of-concept figures and reveals the engineering-gap panel
described in Section~\ref{sec:gap}.

\section{Source Clips}

A card grid of the clip library. Each card shows a procedurally drawn
thumbnail that animates on hover, an active badge, an \emph{Added by you} badge
for uploads, the name, duration and resolution, and either a \emph{Set as
active} action or, for the active clip, a link opening the Console.

An add-clip card opens a modal dialog with a drag-and-drop zone and a file
picker. Non-video files are rejected with an explanatory message. Accepted
files become a new clip immediately and are set active. The uploaded file is
held as an in-memory object URL for the life of the tab; nothing is uploaded,
processed or analysed.

The dialog is rendered through a portal to the document body and implements
standard modal behaviour: focus moves into the panel on open and returns to
the opener on close; Escape closes; a press on the backdrop closes, while a
drag that begins inside the panel does not.

\section{Offline Studio}

The non-real-time mode, presented as three numbered steps with a live output
panel.

\begin{enumerate}
  \item \textbf{Pick a source clip.} A compact picker over the same library.
  \item \textbf{Upload a driving video.} A dropzone using the approved prompt.
    On acceptance the file is previewed, its true duration is read from its
    metadata, and a simulated input-quality check runs for 1.4\,s before
    reporting the four checks as passed. A failure path is available: a file
    whose name contains \emph{bad}, or a demonstration switch, reports
    \emph{Face partially occluded} with the approved recovery message and
    re-locks rendering.
  \item \textbf{Fidelity settings.} An output-resolution selector
    (256, 512, 1024) and a temporal-smoothing slider, with the note that
    offline mode has no latency limit.
\end{enumerate}

Rendering is unlocked only once the input check passes; until then the reason
is stated. The render runs a five-second progress bar through four staged
captions. On completion the output panel reveals the result with a
\emph{Ready} chip, the disclosure overlay if enabled, and a details list
giving the resolution rendered at, the smoothing value, the source clip, the
computed output length, and a latency-budget row reading \emph{None -- offline
quality reference}.

\textbf{Download.} If the active clip was supplied by the user, the original
file is handed back. Otherwise the on-screen stand-in is encoded to a WebM
file whose duration matches the uploaded driving video, clamped to between 2
and 60 seconds; where the browser cannot encode video, a single still frame is
produced instead. The accompanying note states plainly that the file is not
model output.

\section{Devices \& Output}

\begin{itemize}
  \item \textbf{Input device.} A camera selector, cosmetic, with the selection
    reflected on the Console.
  \item \textbf{Output device.} A highlighted card naming
    \emph{NeuroPresence Camera}, carrying the approved guidance for selecting
    it inside a meeting client, and a status chip that reads \emph{Streaming}
    while a session is live and \emph{Ready} otherwise. Supported clients are
    listed with brand marks, alongside the approved limitation note about
    mobile and managed-enterprise clients.
  \item \textbf{How to use in a meeting.} Three steps.
  \item \textbf{Fallback behaviour.} A statement that output falls back to a
    pass-through camera if the pipeline fails or exceeds its latency budget,
    labelled as a stated behaviour rather than a running watchdog.
\end{itemize}

\section{Feasibility evidence}

Described in full in Chapter~\ref{ch:feasibility}.

\section{Settings}

\begin{itemize}
  \item \textbf{Identity.} The enrolled user, the label identifier, the
    enrollment time, and a re-enroll action.
  \item \textbf{Consent policy.} The gate switch, with an explanation that the
    restriction is product policy rather than a user preference, and that the
    finished system targets a true-accept rate of at least 95\% for the
    enrolled user.
  \item \textbf{Disclosure policy.} The watermark switch, with the exact
    overlay wording quoted.
  \item \textbf{Appearance.} Light or dark.
  \item \textbf{About.} Team members with roll numbers, supervisor,
    university, product name and tagline, and the prototype disclosure.
\end{itemize}

Both policy switches are global: flipping either is reflected immediately in
the top bar and on the Console.

%=======================================================================
\chapter{Interaction Flows}
%=======================================================================

Each flow below is completable by clicking, with no dead ends and no console
errors. All six are covered by the automated suite described in
Chapter~\ref{ch:verification}.

\section{Flow A: first run and enrollment}

Launch, arriving on the landing page. \emph{Launch the prototype} opens the
enrollment flow. Step one presents the product; \emph{Get started} advances.
Step two takes a name and runs the two-second simulated analysis, resolving to
an enrolled confirmation. \emph{Continue} advances to step three, where a
source clip is chosen. \emph{Finish} lands on the Console with identity and
clip both set.

A \emph{Skip -- use demo identity} link on step one bypasses the whole flow and
lands on the Console with the demo identity applied, so a demonstration can
start from the centrepiece immediately.

\section{Flow B: a real-time session}

From the Console, \emph{Start Session} runs the warm-up and goes live. The
output frame gains a live indicator, the pipeline strip begins flowing, and
telemetry samples twice per second with the history sparkline filling. The
disclosure watermark can be toggled and the overlay appears and disappears on
the output. The Target and Baseline views can be switched, and the figures
change accordingly. \emph{Stop Session} returns everything to idle, with the
last sample held rather than cleared.

\section{Flow C: demonstrating the safeguard}

With a session live, \emph{Simulate non-enrolled face} flips the gate to
blocked. The running session is torn down, the start control is disabled, a
banner appears, the output frame gains a red veil, and the pipeline strip halts
at the gate with everything downstream dimmed. A countdown runs, and after four
seconds the gate recovers and the start control is enabled again.

The gate can also be switched off entirely, which produces the amber
\emph{disabled} state and greys out the demonstration control.

\section{Flow D: offline mode}

The mode switch or the sidebar opens the Offline Studio. A source clip is
picked, a driving video uploaded, and the input-quality check runs and passes.
Resolution and smoothing are set. \emph{Render} runs the staged progress bar,
and the output panel reveals the result with its details and a download.
\emph{Render again} resets to the pending state.

The failure path is reachable at any point by the demonstration switch or by
uploading a file whose name contains \emph{bad}, and re-locks rendering with an
explanatory message.

\section{Flow E: the meeting-integration story}

Devices \& Output presents the input selector, the named virtual-camera output
with its selection guidance, the supported clients, the desktop-only
limitation, the three-step usage guide, and the fallback behaviour. Starting a
session elsewhere changes the output chip here to \emph{Streaming}.

\section{Flow F: the guided demo}

A single press plays all of the above unattended, narrated, in roughly three
minutes. See Chapter~\ref{ch:tour}.

%=======================================================================
\chapter{Component Catalogue}
%=======================================================================

\section{Interface primitives}

\begin{longtable}{L{3.2cm} L{11.6cm}}
\toprule
\textbf{Component} & \textbf{Description} \\
\midrule
\endhead
\code{Card} & Bordered surface with themed elevation; optional hover lift;
  forwards arbitrary attributes so panels can be addressed by the guided
  demo \\
\code{CardHeader} & Icon, title, subtitle and a right-hand slot \\
\code{Button} & Five variants (primary, secondary, ghost, danger, success),
  three sizes, optional full width, disabled styling \\
\code{Toggle} & Switch with \code{role="switch"}, label, description and three
  tone options \\
\code{Segmented} & Tab group with a spring-animated sliding pill scoped per
  instance \\
\code{Select} & Labelled select with an optional hint \\
\code{ProgressBar} & Determinate bar with full ARIA range semantics \\
\code{Chip} & Six tones for status \\
\code{OverlayChip} & Chip for use over video, dark in both themes \\
\code{SimulatedBadge} & The mandatory honesty marker \\
\code{StatusDot} & Small state dot with optional pulse \\
\code{EmptyState} & Icon, title, body and action \\
\bottomrule
\end{longtable}

\section{Domain components}

\begin{longtable}{L{3.2cm} L{11.6cm}}
\toprule
\textbf{Component} & \textbf{Description} \\
\midrule
\endhead
\code{VideoFrame} & The output surface with its idle, warming, live and blocked
  overlays and the disclosure watermark \\
\code{ClipCanvas} & Procedural clip renderer. Draws a seeded scene -- graded
  background, key light, wall detail, figure, garment and head -- onto a canvas
  with optional idle motion and a talking cadence. Used for every thumbnail and
  preview, so no video file is shipped \\
\code{MetricsPanel} & The telemetry readout \\
\code{Sparkline} & Spline history chart with threshold marker \\
\code{GapMeter} & The measured-versus-target comparison \\
\code{PipelineFlow} & The five-stage pipeline strip \\
\code{ConsentGate} & The gate card and its three states \\
\code{WatermarkCard} & The disclosure toggle and preview \\
\code{SessionControls} & Start, stop, clock and signal path \\
\code{DrivingSignal} & Optional webcam passthrough, display only \\
\code{GuidedTour} & The demo overlay and presenter dock \\
\code{LandingPreview} & The animated product shot on the landing page \\
\code{Avatar} & Initials on a generated gradient; no photograph is ever stored
  because no capture occurs \\
\code{exportOutput} & Canvas-to-WebM encoding and download helpers \\
\code{useAnimatedNumber} & Value interpolation for metric tiles \\
\bottomrule
\end{longtable}

\section{Motion primitives}

\code{Reveal}, \code{RevealGroup} and \code{RevealItem} for scroll-triggered
entrances; \code{MountStagger} and \code{MountItem} for panels already in view
when a screen opens; \code{CountUp} for numerals that count into place;
\code{WordReveal} for headline text rising word by word out of clipping boxes;
\code{TiltCard} for pointer-tracked tilt; \code{AuroraBackdrop} for drifting
colour blooms; \code{GrowBar} for bars that grow into view;
\code{ScrollProgress} for the reading indicator; \code{useScrollSpy} and
\code{scrollToSection} for landing navigation.

Every one of these degrades to an immediate, static render under
\code{prefers-reduced-motion}.

%=======================================================================
\chapter{The Guided Demo}
\label{ch:tour}
%=======================================================================

The rehearsed defense click-path is built into the application. It is launched
from the sidebar or from the landing page, and drives the interface through
eleven steps, operating exactly the same public actions a presenter would click
by hand. Nothing additional is simulated.

\section{Presentation}

A dimming layer covers the viewport with a hole cut over the element under
discussion, tracked by a single animation loop that eases the reported
rectangle toward its target so the mask and its ring stay locked together. A
presenter dock carries the chapter, the title, the line to say and a
supporting note. The dock relocates to the opposite edge when it would
otherwise cover its own subject.

The overlay is click-through everywhere except the dock, so the application
remains drivable by hand at any point without leaving the tour.

\section{Transport}

\begin{center}
\begin{tabular}{ll}
\toprule
\textbf{Control} & \textbf{Action} \\
\midrule
Right / Left arrow & Next or previous step; pauses auto-advance \\
Space & Play or pause \\
Escape & Exit \\
Step dots & Jump to any chapter \\
\bottomrule
\end{tabular}
\end{center}

Auto-advance is on by default so the tour can run unattended; a per-step timer
bar drains over the step's dwell time.

\section{Steps}

\begin{longtable}{L{0.7cm} L{2.3cm} L{3.3cm} L{2.0cm} L{5.4cm}}
\toprule
\textbf{\#} & \textbf{Chapter} & \textbf{Step} & \textbf{Dwell} & \textbf{State the step establishes} \\
\midrule
\endhead
1 & Setup & The product & 7.0\,s & Resets session, render, gate, watermark and
  metrics mode; opens the Console \\
2 & Setup & How it works & 8.5\,s & Spotlights the pipeline strip \\
3 & Live session & Start a session & 6.5\,s & Starts a live session \\
4 & Live session & Live telemetry & 9.0\,s & Ensures a session; target mode \\
5 & The evidence & Measured baseline & 11.0\,s & Switches to baseline; reveals
  the engineering gap \\
6 & Safeguards & Disclosure & 7.5\,s & Enables the watermark \\
7 & Safeguards & The consent gate & 9.0\,s & Blocks the gate; the pipeline
  halts and the session is torn down \\
8 & Two modes & Offline studio & 11.0\,s & Loads a stand-in recording and
  renders it unattended \\
9 & Integration & Virtual camera & 9.0\,s & Opens Devices \& Output \\
10 & Evidence & Feasibility evidence & 10.0\,s & Opens the risk register \\
11 & Evidence & Where this stands & 9.0\,s & Closing statement \\
\bottomrule
\end{longtable}
\noindent\textit{\small The eleven steps. Total unattended runtime is approximately 98
seconds of dwell, or three to four minutes when narrated.}\par\medskip

%=======================================================================
\chapter{Feasibility Evidence}
\label{ch:feasibility}
%=======================================================================

This screen answers the POC-lite criterion of the evaluation rubric directly.
It opens with a coverage strip naming which section satisfies which part of the
criterion.

\begin{longtable}{L{6.6cm} L{8.2cm}}
\toprule
\textbf{Criterion element} & \textbf{Where it is satisfied} \\
\midrule
\endhead
Clickable high-fidelity prototype, no core functionality behind the interface
  & The application as a whole; the scope card on this screen \\
API / data check & Data-path check, this screen \\
Baseline run & Baseline run record, this screen \\
Sample pipeline & The pipeline strip on the Console \\
Risk experiment & Risk register, this screen \\
Full MVP & Not required; deliberately not attempted \\
\bottomrule
\end{longtable}

\section{Data-path check}

There is no backend to probe, so what is checked is the path this build
genuinely depends on. Eight checks run live in the reader's own browser when
the screen opens, and can be re-run on demand. Each reports a real outcome and
a measured duration.

\begin{longtable}{L{4.4cm} L{10.4cm}}
\toprule
\textbf{Check} & \textbf{What it stands behind} \\
\midrule
\endhead
Canvas 2D rendering & Clip stand-ins, the output preview, the render surface.
  Paints a pixel and reads it back \\
Canvas to MediaStream & Encoding the offline result. Opens a capture stream
  and counts tracks \\
WebM video encoder & The download action. Enumerates accepted codec profiles \\
Encode-decode round trip & Reading an uploaded recording's duration. Records a
  real fragment, then decodes its metadata back \\
Object URL lifecycle & Holding uploads in the tab. Creates, reads, revokes,
  and confirms the handle stops resolving \\
Upload type guard & Rejecting non-video uploads. Exercises the exact predicate
  the dialogs use \\
Camera capture API & The optional driving-signal preview. Presence only; no
  permission is requested \\
Reenactment model path & The core functionality deliberately left out.
  Reported as \emph{not implemented}, never as passing \\
\bottomrule
\end{longtable}

\section{Baseline run record}

The environment and the recorded figures, presented as a run record, followed
by the gap panel and a verdict.

\subsection{Engineering gap}
\label{sec:gap}

Both figures are drawn to a common scale so the size of the remaining work is
legible at a glance: 130.1\,ms measured against a 42\,ms target, an 88.1\,ms
reduction, a factor of 3.1. The panel states the hardware, notes that the
figures are measured rather than simulated, and names the optimisation path.

\section{Risk register}

Six risks, each paired with the experiment that settles it.

\begin{longtable}{L{4.1cm} L{1.5cm} L{1.9cm} L{6.9cm}}
\toprule
\textbf{Risk} & \textbf{Severity} & \textbf{Status} & \textbf{Result or response} \\
\midrule
\endhead
Too slow for real-time conversation & High & Run &
  130.1\,ms mean, 7.7\,FPS; 3.1$\times$ above budget. Response: quantisation,
  TensorRT export, lighter warping \\
GPU cannot hold the model & High & Run &
  1.57\,GB inference peak; 6.68\,GB fine-tune fit with roughly 1.3\,GB spare.
  Response: mixed precision if extended \\
Meeting clients refuse a virtual camera & Medium & Partial &
  Desktop clients accept it; mobile and some managed builds block it.
  Response: scope to desktop and state the limit in the interface \\
Identity drifts during a call & High & Planned &
  No result claimed. The metric already has a tile so it is reported from the
  outset \\
Consent gate accepts or rejects wrongly & High & Planned &
  No result claimed. In this build the gate is interface state, not a
  classifier \\
Memory creeps and the process dies & Medium & Planned &
  No result claimed. The fallback it protects is already specified \\
\bottomrule
\end{longtable}

\begin{npnote}{The honesty rule.}
A result is recorded only where the work has actually been done, and every
figure in one traces to the measured run. Rows marked \emph{planned} carry no
number whatsoever, because inventing one would defeat the purpose of the
exercise. This is enforced by an automated test in both directions: no planned
row may display a result, and no completed row may omit one.
\end{npnote}

\section{Scope card}

Two explicit lists: what is implemented and interactive, and what is
deliberately not implemented. This is the criterion's phrase -- an interface
without the core functionality behind it -- stated in the product itself.

%=======================================================================
\chapter{Measured Figures and Targets}
\label{ch:numbers}
%=======================================================================

These are the only non-simulated numbers in the system. They were produced by a
separate benchmark run and are quoted verbatim wherever they appear.

\section{Hardware}

NVIDIA RTX 5050 laptop GPU with 8\,GB of VRAM, running Kubuntu Linux.

\section{Measured proof of concept}

Unoptimized baseline using the published LivePortrait weights, unmodified.

\begin{center}
\begin{tabular}{lr}
\toprule
\textbf{Metric} & \textbf{Value} \\
\midrule
Mean per-frame latency & 130.1\,ms \\
p95 latency & 133.9\,ms \\
p99 latency & 138.7\,ms \\
Throughput & 7.7\,FPS \\
Peak VRAM (inference) & 1.57\,GB \\
Fine-tune fit check (2 of 5 modules, FP32, batch 1) & 6.68\,GB peak \\
\bottomrule
\end{tabular}
\end{center}

\section{Engineering targets}

\begin{center}
\begin{tabular}{lr}
\toprule
\textbf{Attribute} & \textbf{Target} \\
\midrule
Per-frame reenactment compute & $\leq 42$\,ms \\
Throughput & $\geq 24$\,FPS \\
End-to-end latency, capture to virtual camera & $\leq 150$\,ms worst case \\
Peak GPU memory & $\leq 8$\,GB \\
Identity preservation (CSIM) & $\geq 0.80$ \\
Consent-gate true accept, enrolled user & $\geq 95\%$ \\
Endurance & $\geq 30$\,min continuous, no OOM \\
\bottomrule
\end{tabular}
\end{center}

\section{Interpretation}

Memory is comfortable; latency is the whole problem. The gap from 130.1\,ms to
42\,ms is a known optimisation path rather than an open research question, and
closing it is the engineering the project delivers.

%=======================================================================
\chapter{Functional Requirement Traceability}
%=======================================================================

Every functional requirement maps to something visible in the interface.

\begin{longtable}{L{1.3cm} L{5.6cm} L{7.8cm}}
\toprule
\textbf{FR} & \textbf{Requirement} & \textbf{Representation in the prototype} \\
\midrule
\endhead
FR-1 & Capture the live webcam and extract motion &
  The driving-signal thumbnail, the \emph{Driving: Webcam} readout, and the
  first two stages of the pipeline strip \\
FR-2 & Animate a source clip &
  The output preview playing the active source clip; the Reenactment stage \\
FR-3 & Verify likeness before animating &
  The consent gate card and its three states; the gate stage of the pipeline,
  which halts the strip when it blocks \\
FR-4 & Output to a virtual camera selectable by meeting apps &
  Devices \& Output, the named output device, and the final pipeline stage \\
FR-5 & Optional disclosure watermark &
  The watermark card, the overlay on the output frame and on the offline
  render \\
FR-6 & Preserve identity through a session &
  The CSIM telemetry tile with its target threshold \\
FR-7 & Report live latency, throughput and memory &
  The telemetry panel and the history sparkline \\
FR-8 & Offline mode with the same features, driven by a recording &
  The Offline Studio in full \\
\bottomrule
\end{longtable}

%=======================================================================
\chapter{Accessibility and Responsiveness}
%=======================================================================

\section{Keyboard and assistive technology}

\begin{itemize}
  \item Every control is reachable and operable by keyboard, with a visible
    two-pixel focus ring at a two-pixel offset.
  \item Every button carries an accessible name, whether by text, an
    \code{aria-label} or a title.
  \item Every form control is associated with a label.
  \item Switches use \code{role="switch"} with \code{aria-checked}; tab groups
    use \code{role="tab"} with \code{aria-selected}; progress bars carry full
    range semantics; navigation marks the current page with
    \code{aria-current}.
  \item The modal dialog moves focus into itself on open, restores it on
    close, and closes on Escape or a backdrop press.
  \item The guided demo is fully keyboard-operable and exposes its dock as a
    labelled region.
  \item Exactly one first-level heading exists per screen.
\end{itemize}

\section{Responsive behaviour}

The design reference is 1440$\times$900. The sidebar collapses from 252\,px to
a 72\,px icon rail below 1280\,px. Grids reflow at each breakpoint, and the
pipeline strip scrolls horizontally within its own container rather than
forcing the page to scroll. No horizontal page scrolling occurs at any width
down to 480\,px.

\section{Reduced motion}

A global rule collapses animation and transition durations and disables smooth
scrolling. Components additionally branch on the preference so that reveals,
counters, tilts, flow particles and the spotlight resolve immediately to their
final state.

%=======================================================================
\chapter{Verification}
\label{ch:verification}
%=======================================================================

\section{Method}

The prototype is exercised by an automated harness driving a real Chrome
instance over the DevTools Protocol. Tests click actual controls, dispatch real
key events, upload genuinely encoded video files, measure downloaded artifacts
and read computed styles. The full suite is run against both the development
server and the production bundle, because the two differ in ways that have
already hidden a defect once.

\section{Coverage}

\begin{center}
\begin{tabular}{llr}
\toprule
\textbf{Suite} & \textbf{Area} & \textbf{Assertions} \\
\midrule
T1 & Routing, deep links, history, titles & 18 \\
T2 & Console, session lifecycle, telemetry, consent gate & 31 \\
T4 & Enrollment flow, re-enroll, skip & 25 \\
T5 & Offline studio, upload, render, download & 30 \\
T6 & Devices, settings, theme coverage, landing & 50 \\
T7 & Interaction races, responsiveness, reduced motion & 21 \\
T8 & Accessibility and keyboard behaviour & 27 \\
T9 & Guided demo & 32 \\
T10 & Feasibility evidence and the honesty invariant & 37 \\
\midrule
& \textbf{Total} & \textbf{271} \\
\bottomrule
\end{tabular}
\end{center}

All 271 assertions pass against the production build, with no uncaught
exceptions and no console errors.

\section{Notable verifications}

\begin{itemize}
  \item The offline export is measured end to end: a driving video of a known
    length is generated, uploaded, rendered and downloaded, and the resulting
    file's decoded duration is compared against the input.
  \item The guided demo is verified to actually change application state at
    each chapter, not merely to display text.
  \item The risk register's honesty invariant is asserted in both directions.
  \item Rapid navigation, interleaved theme changes and button mashing are
    exercised to confirm no transition can wedge.
\end{itemize}

\section{Defects found and corrected}

\begin{longtable}{L{4.6cm} L{10.2cm}}
\toprule
\textbf{Defect} & \textbf{Cause and resolution} \\
\midrule
\endhead
Every screen blanked after a theme change &
  An exit-then-enter transition deadlocked when an unrelated context update
  re-rendered the subtree, leaving the content region at zero opacity for all
  subsequent navigation. Replaced with a keyed transition that cannot wedge \\
Offline export always three seconds &
  The encoder length was hard-coded. The uploaded recording's true duration is
  now read from its metadata and used, clamped to a sane range \\
Deep links and refresh fell back to the landing page &
  The screen-to-hash effect fired with a stale initial value and overwrote the
  requested route. The initial screen is now seeded from the URL \\
Escape did not close the upload dialog &
  The handler sat on a non-focusable element, so no key event ever reached it.
  Moved to a document-level listener \\
Focus never entered the dialog &
  Keyboard and screen-reader users were stranded. Focus now enters on open and
  returns to the opener on close \\
Modal backdrop 16\,px short &
  A layout utility applied a top margin to the backdrop, and a fixed element
  still honours margins, leaving a strip along the top through which clicks
  fell to the page beneath. The dialog is now portalled to the document body \\
Data checks stuck on \emph{running} &
  A cancellation flag set by the development-mode remount teardown was never
  cleared. Replaced with a run token \\
\bottomrule
\end{longtable}

%=======================================================================
\chapter{Build and Deployment}
%=======================================================================

The production build type-checks the whole source tree and emits static files
to \code{dist/}: a single HTML entry, one stylesheet and one script bundle.
The build uses a relative base path, so the output works from any directory.

Deployment targets any static host. The repository is deployed continuously
from its default branch. Because routing is hash-based, no server rewrite
rules are required.

Requires Node 18 or later.

%=======================================================================
\chapter{Honesty and Non-Implementation}
%=======================================================================

The value of this artifact depends on nobody mistaking it for a working
system. The following measures are structural rather than incidental.

\begin{itemize}
  \item The telemetry panel and the pipeline strip carry a permanent
    \emph{simulated} badge.
  \item The sidebar, the page footer, the About section and the landing page
    each carry the prototype disclosure.
  \item The output preview is described wherever it appears as the source clip
    acting as a stand-in, never as model output.
  \item The consent gate is described as interface state rather than a
    classifier everywhere it appears.
  \item The offline download states that the file is not model output.
  \item The enrollment analysis is described as a timer, and the identifier as
    a label rather than a computed embedding.
  \item The webcam passthrough states that frames are displayed and never
    analysed.
  \item The feasibility screen reports the model path as \emph{not
    implemented} rather than passing it, lists what is deliberately absent, and
    withholds results for experiments not yet run.
  \item All simulated behaviour is confined to a single directory, so the
    boundary can be audited.
\end{itemize}

\begin{npnote}{Statement.}
No machine-learning library appears in the dependency tree. No inference of any
kind is performed. No frame is analysed. No identity is computed or compared.
No virtual camera device is created. No data leaves the browser tab.
\end{npnote}

%=======================================================================
\appendix
%=======================================================================

\chapter{Approved Microcopy}

Fixed strings, reproduced verbatim wherever they appear.

\begin{longtable}{L{3.6cm} L{11.2cm}}
\toprule
\textbf{Context} & \textbf{String} \\
\midrule
\endhead
Tagline & Show up composed in every meeting --- using your own likeness. \\
Enrollment note & NeuroPresence only animates your own consented likeness. \\
Consent verified & Enrolled user --- animation permitted. \\
Consent disabled & Consent gate off --- not recommended. \\
Consent blocked & Likeness does not match the enrolled user --- animation
  blocked. \\
Watermark overlay & Synthetic media --- reenacted. \\
Offline upload prompt & Drop a recording of yourself --- front-facing, face
  clearly visible. \\
Offline note & Offline mode has no latency limit, so higher fidelity settings
  are available. \\
Virtual-camera help & Select 'NeuroPresence Camera' as your camera in Zoom,
  Google Meet, or Microsoft Teams. \\
Limitation note & Desktop apps supported. Mobile and some managed-enterprise
  clients block virtual cameras. \\
Offline failure & We couldn't track the face reliably. Try a clearer,
  front-facing recording. \\
Prototype disclosure & Prototype build. Interface demonstration with simulated
  data; the reenactment pipeline is not implemented in this build. \\
\bottomrule
\end{longtable}

\chapter{Demo Script}

The rehearsed click-path, which the guided demo automates.

\begin{enumerate}\itemsep4pt
  \item \textbf{Open on the Console, session idle.} ``This is NeuroPresence.
    I've already enrolled my identity; my presentable source clip is loaded.''
  \item \textbf{Start Session.} ``Starting a live session --- the system warms
    up, then streams to a virtual camera.''
  \item \textbf{Point at the telemetry panel, in Target mode.} ``This is the
    live telemetry: per-frame latency, frames per second, GPU memory, and
    identity similarity. These are the targets our finished system runs at.''
  \item \textbf{Flip to Baseline.} ``And this is our measured proof of concept
    on the actual RTX 5050 --- 130 ms per frame, 7.7 FPS. That gap, 130 down to
    42 milliseconds, is exactly the engineering our project delivers.''
  \item \textbf{Toggle the disclosure watermark.} ``We can attach a
    synthetic-media disclosure so other participants know the feed is
    reenacted.''
  \item \textbf{Simulate a non-enrolled face.} ``The consent gate only animates
    my own enrolled likeness --- a non-matching face is blocked.''
  \item \textbf{Switch to offline mode and render once.} ``The same feature set
    also runs offline: upload a recording, and with no latency limit we produce
    a higher-fidelity result.''
  \item \textbf{Devices \& Output.} ``Integration is app-agnostic: we present
    as a virtual camera that Zoom, Meet and Teams read as an ordinary webcam
    --- no per-app plugin.''
  \item \textbf{Close.} ``This prototype is the interface; the reenactment
    pipeline is our FYP build. Our proof of concept has already measured
    feasibility on our hardware.''
\end{enumerate}

\begin{npnote}{Note for the presenter.}
Always state that the telemetry in the prototype is simulated, and that
130\,ms is the measured baseline the project intends to reduce to 42\,ms.
Never imply that the prototype performs real reenactment.
\end{npnote}

\chapter{File Inventory}

Thirty-nine source files, approximately 8{,}830 lines of TypeScript and TSX.

\begin{longtable}{L{5.6cm} L{9.2cm}}
\toprule
\textbf{File} & \textbf{Responsibility} \\
\midrule
\endhead
\multicolumn{2}{l}{\textit{Entry and shell}}\\
\code{main.tsx} & Mounts the provider and application \\
\code{App.tsx} & Screen switch, shell composition, footer \\
\code{index.css} & Tokens, base layer, component classes, reduced motion \\
\code{app/routes.ts} & Route table, titles, hash resolution \\
\code{app/Sidebar.tsx} & Primary navigation, demo launcher, identity chip \\
\code{app/TopBar.tsx} & Titles, mode switch, status chips, output name \\
\addlinespace
\multicolumn{2}{l}{\textit{Screens}}\\
\code{screens/Landing.tsx} & Project overview, twelve sections \\
\code{screens/Onboarding.tsx} & Three-step enrollment \\
\code{screens/Console.tsx} & The centrepiece \\
\code{screens/SourceClips.tsx} & Clip library and upload dialog \\
\code{screens/OfflineStudio.tsx} & Offline mode, three steps and output \\
\code{screens/Devices.tsx} & Input, output and integration guidance \\
\code{screens/Feasibility.tsx} & Rubric coverage, checks, run record, risks \\
\code{screens/Settings.tsx} & Identity, policies, appearance, about \\
\addlinespace
\multicolumn{2}{l}{\textit{Components}}\\
\code{components/ui.tsx} & Twelve interface primitives \\
\code{components/motion.tsx} & Fourteen motion primitives and scroll helpers \\
\code{components/VideoFrame.tsx} & Output surface and overlays \\
\code{components/ClipCanvas.tsx} & Procedural clip renderer \\
\code{components/MetricsPanel.tsx} & Telemetry readout \\
\code{components/Sparkline.tsx} & Spline history chart \\
\code{components/GapMeter.tsx} & Measured versus target comparison \\
\code{components/PipelineFlow.tsx} & Five-stage pipeline strip \\
\code{components/ConsentGate.tsx} & Gate card and states \\
\code{components/WatermarkCard.tsx} & Disclosure toggle and preview \\
\code{components/SessionControls.tsx} & Start, stop, clock, signal path \\
\code{components/DrivingSignal.tsx} & Optional webcam passthrough \\
\code{components/GuidedTour.tsx} & Spotlight overlay and presenter dock \\
\code{components/LandingPreview.tsx} & Animated product shot \\
\code{components/Avatar.tsx} & Initials on a generated gradient \\
\code{components/exportOutput.ts} & Canvas encoding and download \\
\code{components/useAnimatedNumber.ts} & Metric value interpolation \\
\addlinespace
\multicolumn{2}{l}{\textit{Simulation}}\\
\code{mock/engine.tsx} & All simulated state and the context provider \\
\code{mock/sampler.ts} & Telemetry random walk \\
\code{mock/constants.ts} & Measured figures, targets, envelopes, microcopy \\
\code{mock/seedData.ts} & Clips, devices, clients, stages, checks \\
\code{mock/tour.ts} & The eleven-step demo script \\
\code{mock/evidence.ts} & Baseline record, risk register, scope \\
\code{mock/dataCheck.ts} & The eight live data-path checks \\
\code{mock/types.ts} & Shared type declarations \\
\bottomrule
\end{longtable}

\chapter{Glossary}

\begin{longtable}{L{3.6cm} L{11.2cm}}
\toprule
\textbf{Term} & \textbf{Meaning} \\
\midrule
\endhead
Source clip & The user's pre-recorded, presentable video, which gets animated.
  Not the live input \\
Driving signal & The live webcam in real-time mode, or an uploaded recording in
  offline mode, supplying motion \\
Reenactment & Transferring the driving motion onto the source clip \\
Consent gate & The safeguard restricting animation to the enrolled user's own
  likeness \\
Disclosure watermark & The optional overlay telling other participants the feed
  is reenacted \\
Virtual camera & A software camera device that meeting applications read as an
  ordinary webcam \\
CSIM & Cosine identity similarity, 0 to 1, higher is better \\
Warm-up & The interval between starting a session and the feed going live \\
Target mode & Telemetry display showing the finished-product operating point \\
Baseline mode & Telemetry display clamped to the measured figures \\
\bottomrule
\end{longtable}

\end{document}
